% Wave-15 R3/20 — slides 400–500 extraction (Nekrasov voice, CG discipline)
% Source: /tmp/automorphic-slides.txt lines 400-500 (Gritsenko–Nikulin programme,
% Lorgat 2020 automorphic-corrections paper). No PDF consulted.

\providecommand{\Lat}{\Lambda}
\providecommand{\II}{\mathrm{II}}
\providecommand{\Aut}{\operatorname{Aut}}
\providecommand{\kch}{\kappa_{\mathrm{ch}}}
\providecommand{\kBKM}{\kappa_{\mathrm{BKM}}}
\providecommand{\kcat}{\kappa_{\mathrm{cat}}}
\providecommand{\PII}{\mathcal{P}_{\II}}
\providecommand{\Cone}{\mathcal{C}}
\providecommand{\Rpos}{\mathbb{R}_{\ge 0}}
\providecommand{\Rspos}{\mathbb{R}_{>0}}

\section*{Slides 400–500: fundamental polyhedron, simple-root Gram, Fourier seed}

The segment sits between Wave~14~Q2 (reflection groups
$W^{(2)}_{I,II}$, Weyl vector introduction) and Wave~14~Q3 (Fourier
expansion of $\chi_5$). Its mathematical content is the explicit
\emph{bridge}: the self-duality of the cone $\Cone(\Lat^{2,1})_+$,
the Gram matrix of the three real simple roots
$\{\delta_1,\delta_2,\delta_3\}\subset \Lat^{2,1}_{\II}$, and the
identification of the Fourier seed $\rho + a$ that opens the
Borcherds product on the next slides.

\paragraph{The $\Aut$-transformation law.} For $w \in W^{(2)}(\Lat^{2,1}_{I})$,
$a \in \Aut(\mathcal{P}_{I,\mathrm{prim}})$,
\[
  \Delta_5(w\cdot az) \;=\; \det(a)\,\Delta_5(z),
\]
and for $w \in W^{(2)}(\Lat^{2,1}_{\II})$,
$a \in \Aut(\PII,\mathrm{prim})$,
\[
  \Delta_5(w\cdot az) \;=\; \det(w)\,\Delta_5(z).
\]
The symmetry interchange $\det(a) \leftrightarrow \det(w)$ between
the two lattice frames is the shadow of the birational equivalence
between the $I$- and $\II$-primitive fundamental domains: one half
of the stabiliser absorbs the sign into the $\Aut$-factor, the other
half into the reflection-group factor.

\paragraph{The fundamental polyhedron as self-dual slice.} Let
$\delta_1,\delta_2,\delta_3$ span the positive orthant
\[
  \Delta(\Lat^{2,1}_{\II})_+ \;=\; \Rpos\delta_1 \oplus \Rpos\delta_2 \oplus \Rpos\delta_3,
\]
with dual cone
\[
  \Delta(\Lat^{2,1}_{\II})^{*}_{+} \;=\; \{\, x \in \Lat^{2,1}\otimes\mathbb{R} \;:\; (x,\delta_i) \le 0 \,\}.
\]
The finite-volume condition $\PII \subset \Cone(\Lat^{2,1})_+ / \Rspos$
is equivalent to the double inclusion
\[
  \Delta(\Lat^{2,1}_{\II})^{*}_{+}
  \;\subset\;
  \Cone(\Lat^{2,1})_+
  \;=\;
  \Cone(\Lat^{2,1})^{*}_{+}
  \;\subset\;
  \Delta(\Lat^{2,1}_{\II})_+.
\]
The middle equality is the \emph{self-duality of the signature-$(2,1)$
light cone}: a Lorentzian fact which on lines 422–427 is deployed
twice, once under each frame ($I$ and $\II$). Self-duality is not
decoration; it is what forces the primitive Weyl chamber to sit
inside its own polar.

\paragraph{The Gram matrix and the Weyl vector.} The three real
simple roots on $\Lat^{2,1}_{\II} = \mathbb{Z}\delta_1\oplus\mathbb{Z}\delta_2\oplus\mathbb{Z}\delta_3$
have Gram
\[
  (\delta_i,\delta_j) \;=\;
  \begin{pmatrix} 2 & -2 & -2 \\ -2 & -2 & 2 \\ -2 & 2 & -2 \end{pmatrix}.
\]
The signature $(2,1)$ shows in the spectrum; the hyperbolicity
is visible without diagonalisation from $\det = +16$ and the
off-diagonal $\pm 2$ pattern. The lattice Weyl-vector equations
\[
  (\rho,\delta_i) \;=\; -\tfrac{1}{2}(\delta_i,\delta_i) \;=\; -1
\]
admit the solution
\[
  \rho \;=\; \tfrac{1}{2}\delta_1 + \tfrac{1}{2}\delta_2 + \tfrac{1}{2}\delta_3
       \;=\; f_2 - \tfrac{1}{2} f_3 + f_{-2},
\]
where $\{f_{-2},f_3,f_2\}$ is the rescaled hyperbolic frame with
$(\Lat^{2,1}_{\II})^* = \tfrac{1}{2}\mathbb{Z} f_2 \oplus \tfrac{1}{2}\mathbb{Z} f_3 \oplus \tfrac{1}{2}\mathbb{Z} f_{-2}
 = \tfrac{1}{2}\Lat^{2,1}$. The identification
$\rho \in \Delta(\Lat^{2,1}_{\II})^{*}_{+}$ is the statement that the
Weyl vector lies in the polar of the simple-root cone — the
\emph{inside} of the fundamental polyhedron, not its wall.

Cross-check against Lorgat 2020 (memory: Gram-matrix normal form with
$+2$ on the diagonal and $\pm 2$ off-diagonal, signature $(2,1)$,
$\rho = \tfrac{1}{2}(\delta_1+\delta_2+\delta_3)$): slides and paper
agree; the slide's matrix presentation uses the permutation convention
in which the off-diagonal pattern is $(-,-,+,-,+,-)$ rather than all
$-2$, encoding the choice of ordered basis consistent with Gritsenko–Nikulin
1998.

\paragraph{The Fourier seed.} With $z = z_1 f_{-2} + z_2 f_3 + z_3 f_2
\in \Lat^{2,1}\otimes\mathbb{R} + i\Cone(\Lat^{2,1})_+$ and
integers $n,m > 0$, $l$ with $n \equiv m \equiv l \equiv 1 \pmod 2$
and $4nm - l^2 > 0$, the Fourier coefficient is absorbed into the
$\rho$-shift
\[
  \tfrac{1}{64}\,f(n,l,m)\,e^{\pi i(nz_1+lz_2+mz_3)}
  \;=\; m(a)\,e^{-\pi i(\rho+a,z)},
\]
with
\[
  a \;=\; (n-1)f_2 - (l-1)\tfrac{1}{2}f_3 + (m-1)f_{-2}
       \;\in\; (\Lat^{2,1})^{*} = \tfrac{1}{2}\Lat^{2,1}_{\II},
\]
and $m(a) = -\tfrac{1}{64} f(n,l,m) \in \mathbb{Z}$, $m(0) = -1$.

This is the \emph{single identity} that converts a naive Jacobi
Fourier sum into a product supported on the positive cone: every
$(n,l,m)$ congruence-triple produces a lattice vector $\rho + a$ in
$\Cone(\Lat^{2,1})_+$, and the integrality $m(a) \in \mathbb{Z}$ is the
denominator-formula seed of the Borcherds product for $\Delta_5$. The
normalisation $m(0) = -1$ fixes the leading Weyl-vector term
$-e^{-\pi i(\rho,z)}$ that the next slide-block (Wave~14~Q3, 500–600)
identifies with the Borcherds denominator's $-q^{\rho}$.

\paragraph{What this segment uniquely contributes.} Waves~11–14
named the reflection groups, fixed the lattice conventions, and
opened the Fourier expansion. Lines 400–500 are what makes those
three threads a single computation: (i) the \emph{self-duality}
identification that turns the abstract finite-volume condition into
the computable double inclusion, (ii) the \emph{explicit Gram matrix}
with its off-diagonal $\pm 2$ pattern from which the Weyl vector is
read off directly rather than guessed, and (iii) the
\emph{$\rho$-shift change of variables} $f(n,l,m) \mapsto m(a)$ that
is the precise input to the Borcherds lift
$\mathrm{Lift}(\phi_{0,1}) = \Delta_5$ of the Lorgat 2020 paper.

\paragraph{Frontier.} Within the Vol~III programme the segment
witnesses no genuinely novel structure beyond the established
$\kBKM(\Phi_1) = c_1(0)/2 = 5$ identity and the Wave~14
adjudication ledger (Monster rank 2, Bruinier $c_3 = -8$, admissible
Heegner scope). It is the \emph{foundation stone} on which those
verified numerics rest: the Gram matrix fixes $\rho$, $\rho$ fixes the
Fourier seed, the Fourier seed fixes the Borcherds product, and the
Borcherds product fixes $\kBKM = 5$. The Nekrasov reading is that the
hyperbolic Gram $(\delta_i,\delta_j)$ is the classical shadow of the
K3 $\times$ E chiral algebra's three-dimensional imaginary cone;
quantisation of this cone is precisely what $\Phi$ delivers on the
$N=1$ row.
